% =============================================================================
% Project : THE_PHILOSOPHY_OF_VIABILITY
% Book    : Book I — Ontology of Viability
% Book ID : VS.B01
%
% File    : 2_MAINMATTER/PART_I_IDENTITY/chapter_01_identity/sections/section_08_identity_and_signature.tex
%
% Section / Chapter:
%   Identity and Signature — Chapter 01: Identity
%
% Purpose :
%   To examine the relation between ontological identity and signature,
%   distinguish signature from determination, invariance, numerical identity,
%   individuation, and manifestation, and clarify the unresolved ontological
%   status of signature.
%
% =============================================================================

\section{Identity and Signature}
\label{sec:identity-and-signature}

\subsection{The Ontological Problem of Signature}

The relation between identity and signature raises a distinct
problem within the ontology of identity: whether an organized
configuration of determinations can characterize an existent
without thereby constituting, or exhausting, its identity.

An existent may involve multiple determinations and relations
that form an organized configuration. Some of these
determinations may be materially, structurally, or relationally
connected. Their configuration may provide a significant basis
for characterizing the existent in a particular respect. The
term \emph{signature} will be used for such an organized
configuration when considered in its characterizing role.

This use of the term does not yet settle the ontological status
of signature. In particular, it does not determine whether a
signature is merely an analytical configuration of actual
determinations, whether it corresponds to a distinct
ontological organization, or whether its status differs across
different kinds of existents.

The central question is therefore not simply whether a
signature can characterize an existent. It is whether the
characterizing role of a signature warrants identifying it
with the existent's identity, or assigning it a constitutive
role in that identity.

The distinction between identity and signature must be
maintained while these possibilities are examined.

Identity concerns what makes an existent the existent that
it is, without presupposing that its identity is reducible
to any single determination or configuration. Signature,
by contrast, concerns an organized configuration of
determinations through which an existent may be characterized
in a particular respect.

The two are related because a signature may involve
determinations that are relevant to the constitution,
character, or differentiation of an existent. Their relation,
however, does not establish their equivalence. Nor does the
fact that a signature characterizes an existent establish
that it provides a complete account of what that existent
is.

The present section therefore rejects the identification
of identity with signature while leaving open the more
specific question of whether, and under what conditions,
signature may participate in the ontological constitution
of identity.

\subsection{Identity Is Not Reducible to Signature}

Identity cannot be defined simply as the signature of an
existent.

Such a definition would make the identity of an existent
equivalent to a configuration of determinations selected
and organized as its signature. This would conflate the
ontological question of what the existent is with the
question of which determinations are brought together
to characterize it.

Several difficulties follow from this reduction.

First, a signature may be selective. It may include some
determinations while omitting others. If identity were
equivalent to signature, every omitted determination
would risk being treated as irrelevant to identity, even
where that determination might be ontologically
constitutive.

Second, signatures may be specified at different levels
of comprehensiveness or under different respects of
consideration. The same existent may consequently admit
different characterizations without there being sufficient
grounds to conclude that its identity has changed.

Third, determinations included in a signature may vary.
A change in one or more of them does not, by itself,
establish a change in numerical identity or in the
ontological constitution of the existent.

Fourth, the same configuration may characterize more
than one numerically distinct existent. A signature's
descriptive or structural distinctiveness within one
domain cannot therefore be assumed to provide a
universal criterion of individuation or numerical identity.

These difficulties do not show that signatures are
ontologically insignificant. They show that their
significance cannot be established by identifying them
with identity in advance.

A signature may be relevant to the constitution or
characterization of an existent. Whether that relevance
is constitutive, conditional, or merely characterizing
must be examined independently.

\subsection{Signature as an Organized Configuration}

A signature is not adequately understood as an arbitrary
list of properties. Its characterizing role depends on
the organization of the determinations involved.

A configuration may be organized through relations among
its determinations, through their arrangement, or through
the way in which they jointly characterize an existent.
The significance of the configuration may therefore
depend not only on which determinations are included,
but also on how they are related.

This does not imply that every signature possesses a
single organizing principle, that its components form
a unified whole, or that the organization is necessarily
intrinsic to the existent. These are further questions
that cannot be settled by the concept of signature alone.

Nor does the organization of a configuration establish
the ontological priority of its components. A determination
may be relevant to a signature without being ontologically
prior to the other determinations involved. Similarly,
the coherence of a configuration does not establish that
all its components are constitutive of identity.

The term signature thus identifies an organized
configuration considered in relation to the
characterization of an existent. It does not, by itself,
determine whether that organization constitutes an
ontological unity, an analytical grouping, or a
combination of both.

This distinction is necessary because an existent's
actual determinations and the configuration through
which some of them are considered are not automatically
the same thing. The determinations may belong to the
existent independently of whether they are grouped
together as a signature.

\subsection{Signature and Determination}

Signature and determination are related but non-equivalent.

Determination concerns the ontological condition through
which an existent is determinate. Signature concerns an
organized configuration involving determinations that
may be relevant to characterizing that existent.

A signature therefore does not constitute the totality
of an existent's determinations merely by bringing
some of them together. Its scope may be selective,
respect-dependent, or limited to a particular
configuration under consideration.

The omission of a determination from a signature does
not establish that the determination is absent from
the existent. Nor does it establish that the omitted
determination lacks ontological significance.

Conversely, the inclusion of a determination in a
signature does not establish that it is constitutive
of identity. Its inclusion establishes, at most, its
relevance to the configuration as specified. The
nature and extent of that relevance require separate
examination.

It is also necessary to distinguish the actual
determinations of an existent from the criteria through
which a signature is specified. A criterion of inclusion
may determine which features are considered together,
but it does not thereby determine which features
actually belong to the existent.

Signature is consequently neither an alternative
name for determination nor an exhaustive inventory
of determination. It is a configuration involving
determinations whose ontological relations and
significance cannot be inferred solely from their
participation in that configuration.

\subsection{The Ontological Status of Signature}

The distinction between identity and signature does
not, by itself, establish what kind of ontological
status signature possesses.

At least three possibilities must be distinguished.

First, a signature may be an analytical configuration:
a grouping of actual determinations considered together
for the purpose of characterizing an existent. On this
account, the configuration as a signature need not
constitute an additional ontological feature of the
existent, even though the determinations grouped
together are real.

Second, a signature may correspond to an organization
that belongs to the existent independently of its
description. In this case, the signature would not
merely collect determinations analytically; it would
also refer to an organization grounded in the existent
itself. Whether such an organization is constitutive
of identity would remain a further question.

Third, signature may have a conditional or domain-specific
ontological significance. Certain configurations may
participate in the constitution of particular kinds
of existents or in particular respects of their
identity, without signature having one universal
ontological function.

These possibilities should not be treated as an
exhaustive classification. They identify distinctions
that must be investigated before a general ontological
account of signature can be established.

In particular, the claim that a signature corresponds
to an actual organization does not entail that it
constitutes identity. Likewise, the claim that a
signature is analytically specified does not entail
that the determinations involved are merely descriptive
or lack ontological significance.

The ontological status of signature therefore remains
an open problem. The present analysis establishes
the non-equivalence of identity and signature, not
a definitive theory of signature's independent
ontological status.

\subsection{Signature and Invariance}

Signature may be related to invariance when some
determinations or relations within a configuration
remain preserved under specified transformations.

Invariance concerns the preservation of a specified
feature, relation, or configuration under a particular
class of transformations and with respect to a
specified domain. Signature concerns an organized
configuration of determinations considered in its
characterizing role.

A signature may contain invariant and variable
determinations. Its configuration may remain stable
under some transformations while changing under
others. Whether such changes affect the relevance
or organization of the signature depends on the
configuration and the transformations being considered.

Conversely, an invariant may be shared by multiple
numerically distinct existents. The preservation of
a feature under transformation therefore does not
necessarily distinguish one existent from another.

Neither the stability of a signature nor the
presence of invariants within it establishes that
the signature is identical to identity. Nor does
the absence of an identified invariant establish
that an existent lacks identity.

An earlier formulation of identity in terms of
invariance under allowable transformations cannot
therefore be transferred into the present account
without additional argument. Such a formulation
would require an independently justified account
of the relevant transformations, the domain of
application, and the respect in which preservation
is being considered.

Invariance may provide a condition for examining
the stability of particular signatures. It cannot
be assumed to provide a universal definition of
signature or identity.

\subsection{Signature and Numerical Identity}

Signature cannot be treated as a sufficient criterion
of numerical identity.

Two numerically distinct existents may possess
identical or indistinguishable signatures within
a specified domain and respect of comparison.
The sameness of their signatures therefore does
not establish that they are one and the same existent.

Conversely, the signature associated with one
existent may change as some of its determinations
or relations change. Such a change does not, by
itself, establish that the existent has ceased
to exist or that another existent has replaced it.

These possibilities show that the comparison of
signatures cannot independently settle the question
of numerical identity. Signature similarity does
not establish numerical identity, and signature
difference does not establish numerical non-identity.

This conclusion must be distinguished from the
question of persistence. Numerical identity concerns
whether different designations or considerations
refer to one and the same existent. Persistence
concerns whether an existent continues to exist
through relevant changes or conditions. A signature
may be relevant to an investigation of persistence,
but its stability is not, without further argument,
a sufficient criterion of persistence.

The ontological role of signature in the continuity
of an existent must therefore be investigated
without presupposing that the preservation of a
signature is either necessary or sufficient for
the persistence of that existent.

\subsection{Signature and Individuation}

Signature may characterize an individual existent
in relation to other existents. It does not thereby
explain how that existent is constituted as an
individual.

Individuation concerns the ontological constitution
of an existent as an individual. Signature concerns
an organized configuration of determinations through
which that individual may be characterized.

A signature may be distinctive within a particular
domain or respect of comparison. Such distinctiveness
may assist in differentiating an existent from
others under those conditions. It does not,
however, establish that the signature provides
the ontological ground of individuation.

Likewise, the absence of a unique signature does
not establish the absence of individuality. Two
numerically distinct individuals may share a
signature, or may be indistinguishable with respect
to the determinations included in a particular
configuration.

It follows that individuality must not be made
dependent upon the possession of an absolutely
unique signature. The possibility of shared
signatures does not eliminate numerical distinction,
just as a distinctive signature does not, by itself,
explain the constitution of an individual.

Signature may therefore contribute to the
characterization of individuality without being
identified with individuation or treated as its
universal ontological ground.

\subsection{Signature and Manifestation}

Signature and manifestation concern different
aspects of an existent.

Manifestation concerns the ways in which an existent
or its determinations become manifest under
particular ontological conditions. Signature
concerns an organized configuration of determinations
through which an existent may be characterized.

A particular manifestation may involve some
determinations associated with a signature without
exhausting that signature. Different manifestations
may also involve different subsets or arrangements
of determinations.

Such variation does not automatically establish
a change in identity. Nor does the fact that a
signature characterizes an existent establish
that the signature is itself a manifestation.

Three questions must therefore remain distinct:

\begin{itemize}
    \item What constitutes the ontological identity
    of an existent?

    \item Which organized configuration of
    determinations characterizes the existent
    in a specified respect?

    \item Under what ontological conditions do
    particular determinations become manifest?
\end{itemize}

These questions may constrain one another, but
their answers cannot be presumed to coincide.

The relation between signature and manifestation
also remains open in a further respect: a signature
may be associated with multiple manifestations,
while a manifestation may involve determinations
that are not included in a particular signature.
The precise relation depends on the organization
and scope of the configuration under consideration.

The present analysis consequently rejects both
the reduction of identity to signature and the
reduction of signature to manifestation.

\subsection{The Possibility of Multiple Signatures}

The distinction between identity and signature
allows for the possibility that one existent may
be characterized through more than one signature.

Different configurations may organize different
determinations or relations, or characterize
different respects of the same existent. Their
differences need not imply that the existent
possesses multiple identities.

Nor should the possibility of multiple signatures
be taken to establish that all such configurations
are equally comprehensive or ontologically
equivalent. One configuration may include
determinations omitted from another; two
configurations may overlap; or they may concern
different domains and respects.

These possibilities require distinctions among
the relations that may hold between signatures.
Overlap does not entail identity of configurations.
Partial inclusion does not establish ontological
priority. Difference between configurations
does not, by itself, establish difference between
the identities of the existents they characterize.

The possibility of multiple signatures must also
be distinguished from the question of whether
one existent can involve multiple identities.
The former concerns the plurality of configurations
through which an existent may be characterized.
The latter concerns the ontological constitution
of identity itself.

The possibility of multiple signatures is therefore
compatible with leaving the multiplicity of identity
unresolved. It neither establishes nor excludes
that possibility.

\subsection{Signature, Description, and Ontological Reality}

The characterization of an existent through a
signature must not be confused with the production
of the existent or its determinations through
description.

An existent does not acquire its identity because
a signature has been specified for it. Its actual
determinations do not depend for their existence
upon whether they have been collected into a
configuration or represented in a description.

A description may select, order, and represent
determinations associated with an existent. The
resulting representation may be more or less
comprehensive, accurate, or adequate to the
configuration it purports to characterize. These
features concern the description and its relation
to what it describes; they do not establish that
the existent itself is constituted by the
description.

The distinction between an actual configuration
and its description is therefore essential.
A description of a signature is not identical
to the signature, if the latter corresponds to
an organization grounded in the existent. If
signature is instead used analytically for a
selected configuration, the description still
does not create the determinations it brings
together.

The absence of a complete description does not
establish the absence of identity, determination,
or an actual organization. Nor does the
specification of a complete description establish
that the identity of the existent has been
exhaustively accounted for.

This preserves the ontological independence
of the existent from its characterization while
leaving open the question of whether signature
refers to an organization that possesses an
ontological status beyond its analytical
specification.

\subsection{The Limits of Signature as an Account of Identity}

The analysis establishes several constraints
on the use of signature within an ontology
of identity.

First, signature is not synonymous with identity.
An organized configuration of determinations
cannot be assumed to exhaust what makes an
existent the existent that it is.

Second, signature is not identical to the totality
of an existent's determinations. A configuration
may be selective without implying that omitted
determinations are absent or ontologically
irrelevant.

Third, the inclusion of a determination in a
signature does not by itself establish that
the determination is constitutive of identity.
Its constitutive significance requires an
independent argument.

Fourth, the organization of a signature does
not establish that it constitutes an ontological
unity or possesses an independent ontological
status.

Fifth, signature must not be reduced to invariance.
A signature may include both variable and
invariant determinations, and an invariant may
be shared by distinct existents.

Sixth, the sameness or difference of signatures
does not independently establish numerical
identity or numerical non-identity.

Seventh, a distinctive signature does not
establish the ontological ground of individuation,
and the absence of a unique signature does
not eliminate individuality.

Eighth, a change in signature does not
automatically establish a change in identity,
the end of an existent, or its replacement.

Ninth, multiple signatures associated with
one existent do not, by themselves, imply
multiple identities.

Finally, neither the analytical specification
of a signature nor its descriptive usefulness
settles its ontological status.

These constraints do not render signature
superfluous. They establish the conditions
under which its significance can be examined
without conflating characterization with
ontological constitution.

Signature remains relevant as a way of
considering organized configurations of
determinations. Whether particular configurations
also participate in the constitution of identity
is a further question, to be settled through
specific ontological arguments rather than
through the concept of signature alone.

\subsection{Conceptual Structure of Identity and Signature}

The distinctions developed in this section
can be summarized by separating identity,
determination, signature, invariance, numerical
identity, individuation, manifestation, and
description according to the questions each
concept addresses.

The following table records these distinctions.
It is an analytical guide, not a final taxonomy
of independent ontological entities or domains.

% =============================================================================
% Project : THE_PHILOSOPHY_OF_VIABILITY
% Book    : Book I — Ontology of Viability
% Book ID : VS.B01
%
% File    : 4_ASSETS/tables/VS.B01.P01.C01.S08.SS15.P001.TL01.identity_signature_conceptual_distinctions.tex
%
% Section / Chapter:
%   Identity and Signature — Chapter 01: Identity
%
% Asset:
%   TL01 — Identity and Signature: Conceptual Distinctions
%
% Purpose :
%   To distinguish identity from signature and related concepts,
%   clarifying their different analytical functions and preventing
%   unsupported ontological reductions.
%
% =============================================================================

\begin{table}[htbp]
    \centering
    \caption{Conceptual Distinctions Between Identity and Signature}
    \label{tab:identity-signature-conceptual-distinctions}

    \small
    \renewcommand{\arraystretch}{1.35}

    \begin{tabular}{p{0.19\textwidth} p{0.35\textwidth} p{0.36\textwidth}}
        \hline
        \textbf{Concept}
        &
        \textbf{Central Concern}
        &
        \textbf{Distinction from Signature}
        \\
        \hline

        Identity
        &
        What makes an existent the existent that it is.
        &
        Cannot be reduced to an organized configuration of determinations.
        \\

        Determination
        &
        The ontological condition through which an existent is determinate.
        &
        Signature involves selected and organized determinations but is not identical to determination as such.
        \\

        Signature
        &
        An organized configuration of determinations through which an existent may be characterized in a particular respect.
        &
        Its ontological status and possible constitutive role require independent examination.
        \\

        Invariance
        &
        The preservation of a specified feature, relation, or configuration under particular transformations.
        &
        May concern the stability of some signature components but does not define signature or identity.
        \\

        Numerical Identity
        &
        Whether different designations or considerations refer to one and the same existent.
        &
        Neither identical nor reducible to the sameness of signatures.
        \\

        Individuation
        &
        The ontological constitution of an existent as an individual.
        &
        A distinctive signature may characterize an individual without establishing its ontological ground.
        \\

        Manifestation
        &
        The ways in which an existent or its determinations become manifest under particular ontological conditions.
        &
        A manifestation may involve some signature determinations without exhausting the signature.
        \\

        Description
        &
        The representation or articulation of an existent or a configuration of its determinations.
        &
        A description of a signature must not be confused with the actual determinations or organization it represents.
        \\

        \hline
    \end{tabular}

    \vspace{0.5em}

    \begin{minipage}{0.92\textwidth}
        \footnotesize
        \textit{Note:} This table distinguishes analytical functions
        and ontological questions. It does not establish that every
        listed concept corresponds to an independent ontological
        entity or that the ontological status of signature has been
        definitively settled.
    \end{minipage}
\end{table}

% =============================================================================
% END OF TABLE TL01
% =============================================================================

The following diagram represents the principal
relations and non-equivalences examined in
this section. It distinguishes the existent
and its determinations from the configuration
considered as a signature and from the
description through which that configuration
may be represented.

% =============================================================================
% Project : THE_PHILOSOPHY_OF_VIABILITY
% Book    : Book I — Ontology of Viability
% Book ID : VS.B01
%
% File    : 4_ASSETS/diagrams/VS.B01.P01.C01.S08.SS15.P001.DG01.identity_signature_conceptual_relations.tex
%
% Section / Chapter:
%   Identity and Signature — Chapter 01: Identity
%
% Asset:
%   DG01 — Identity and Signature: Conceptual Relations
%
% Purpose :
%   To represent the conceptual distinction between an existent,
%   its determinations, an organized configuration considered as
%   a signature, and the description through which that configuration
%   may be represented.
%
% Technical Requirements (preamble):
%   \usepackage{tikz}
%   \usetikzlibrary{arrows.meta, positioning}
%
% =============================================================================

\begin{figure}[htbp]
    \centering
    \hyphenpenalty=10000 \exhyphenpenalty=10000\relax  % no hyphenation inside diagram labels

    \begin{tikzpicture}[
        font=\small,
        node distance=0.9cm and 0.7cm,
        % --- base box ---
        box/.style={
            draw, rounded corners=3pt, align=center,
            text width=3.4cm, minimum height=1.6cm,
            inner sep=5pt, line width=0.7pt
        },
        % --- node roles ---
        onto/.style    ={box, fill=blue!8, draw=blue!55!black},
        analytic/.style={box, dashed, fill=gray!10, draw=black!60},
        % --- arrows ---
        directed/.style={-{Stealth[length=2.5mm]}, thick},
        assoc/.style   ={{Stealth[length=2.2mm]}-{Stealth[length=2.2mm]},
                         thick, dashed}
    ]

    % ---------------------------------------------------------------
    % Grid: three columns, centered on Determinations
    % ---------------------------------------------------------------
    \node[onto] (determinations) {
        \textbf{Determinations}\\[2pt]
        Actual determinations of the existent
    };
    \node[onto, left=of determinations] (identity) {
        \textbf{Identity}\\[2pt]
        What makes the existent the existent that it is
    };
    \node[onto, right=of determinations] (manifestation) {
        \textbf{Manifestation}\\[2pt]
        Manifestation of an existent or its determinations
    };
    \node[onto, above=of determinations] (existent) {
        \textbf{Existent}\\[2pt]
        Ontological actuality
    };

    % ---------------------------------------------------------------
    % Analytical column (below Determinations)
    % ---------------------------------------------------------------
    \node[analytic, below=of determinations] (signature) {
        \textbf{Signature}\\[2pt]
        Organized configuration of determinations
    };
    \node[analytic, below=of signature] (description) {
        \textbf{Description}\\[2pt]
        Representation or articulation of a signature
    };

    % Numerical Identity: same column as Identity, same row as Signature
    \node[onto] (numerical) at (identity |- signature) {
        \textbf{Numerical Identity}\\[2pt]
        Reference to one and the same existent
    };

    % ---------------------------------------------------------------
    % Relations
    % ---------------------------------------------------------------
    \draw[directed] (existent)        -- (determinations);
    \draw[directed] (determinations)  -- (signature);

    \draw[assoc] (existent.west) -| (identity.north);
    \draw[assoc] (determinations)     -- (manifestation);
    \draw[assoc] (signature)          -- (description);
    \draw[assoc] (identity)           -- (numerical);

    % ---------------------------------------------------------------
    % Legend (bottom right, in the empty cell)
    % ---------------------------------------------------------------
    \node[anchor=center, draw=black!30, rounded corners=2pt,
          text width=3.4cm, align=left, font=\scriptsize, inner sep=5pt]
        (legend) at (manifestation |- description) {
        \textbf{Legend}\\[3pt]
        \tikz[baseline=-0.5ex]\draw[rounded corners=1pt, thick, fill=blue!8,
            draw=blue!55!black] (0,-0.1) rectangle (0.5,0.1);\;
        ontological concept\\[2pt]
        \tikz[baseline=-0.5ex]\draw[rounded corners=1pt, thick, dashed,
            fill=gray!10, draw=black!60] (0,-0.1) rectangle (0.5,0.1);\;
        analytical notion\\[2pt]
        \tikz[baseline=-0.5ex]\draw[directed] (0,0) -- (0.5,0);\;
        directed relation\\[2pt]
        \tikz[baseline=-0.5ex]\draw[assoc] (0,0) -- (0.5,0);\;
        association
    };

    % ---------------------------------------------------------------
    % Explanatory note
    % ---------------------------------------------------------------
    \node[below=0.8cm of description, text width=11cm, align=center,
          font=\footnotesize\itshape] (note) {
        Conceptual associations do not establish equivalence,
        causal production, or ontological reduction.
        The constitutive status of signature remains open.
    };

    \end{tikzpicture}

    \caption{Conceptual Relations Between Identity, Determination,
    Signature, and Description}
    \label{fig:identity-signature-conceptual-relations}
\end{figure}

% =============================================================================
% END OF DIAGRAM DG01
% =============================================================================

The diagram is conceptual rather than causal,
temporal, or hierarchical. It does not imply
that identity produces signature, that signature
produces invariance, or that manifestation
follows from signature. Nor does it establish
that signature is an independent ontological
layer.

In particular, the distinction between an
existent's actual determinations, their
organization, and the analytical specification
of a signature must not be mistaken for a
settled account of the ontological status of
that organization. The diagram preserves
the questions that remain open rather than
resolving them through visual arrangement.

\subsection{Concluding Formulation}

Identity and signature are related but
non-equivalent.

Identity concerns what makes an existent
the existent that it is. Signature concerns
an organized configuration of determinations
through which an existent may be characterized
in a particular respect.

A signature may characterize determinations
that are relevant to an existent's constitution.
It may also correspond, in some cases, to an
organization grounded in the existent itself.
Neither possibility establishes that signature
is identical to identity or that it exhausts
the ontological constitution of the existent.

The central principle of this section is
therefore:

\begin{quote}
The signature of an existent is not identical
to its ontological identity. It is an organized
configuration of determinations through which
the existent may be characterized, and whose
ontological status and possible constitutive
role must be established independently.
\end{quote}

This principle preserves the distinction
between an existent, its identity, the
determinations associated with it, and the
configurations through which those
determinations may be considered.

It also leaves open a substantive ontological
question: whether signature is exclusively
a mode of analytical characterization or
whether particular signatures may correspond
to, or participate in, real organizations
within the constitution of existents.

The non-reducibility of identity to signature
does not settle that question. It establishes
the conceptual and ontological distinctions
that any adequate answer must respect.

% =============================================================================
% END OF SECTION 08 OF CHAPTER 01
% =============================================================================
